\documentclass[11pt]{article}
\usepackage[margin=1in]{geometry}
\usepackage[T1]{fontenc}
\usepackage[utf8]{inputenc}
\usepackage{lmodern}
\usepackage{hyperref}
\usepackage{enumitem}
\usepackage{longtable}
\usepackage{tabularx}
\usepackage{needspace}
\setlist[itemize]{leftmargin=1.25em}
\setlist[enumerate]{leftmargin=1.5em}

\title{MarketOS User Manual\\\large CSI43C9 Farmers Market Application}
\author{}
\date{\today}

\begin{document}
\maketitle
\tableofcontents
\newpage

\section{What This App Is}
MarketOS is a farmers market management app. It helps staff keep everything in one place:
\begin{itemize}
  \item vendor information,
  \item weekly market-day transactions,
  \item program token activity (SNAP, DUFB, WDFM, vouchers),
  \item market goals,
  \item financial reports,
  \item and admin-level setup.
\end{itemize}

If you are brand new, this manual walks you through each screen in plain language.

\section{Quick Start (For Non-Technical Users)}
\subsection{Open the App}
Depending on how your team runs it, you may use:
\begin{itemize}
  \item \textbf{Desktop app} (preferred for day-to-day use), or
  \item \textbf{Browser app} (development/testing setup).
\end{itemize}

If your team already installed it, open \textbf{MarketOS} and skip to Section \ref{sec:main-navigation}.

\subsection{If You Need to Run It Locally}
Ask your technical lead for help with first-time setup. The short version:
\begin{enumerate}
  \item Start backend server (in \texttt{backend}): \texttt{./gradlew bootRun}
  \item Start frontend (in \texttt{frontend}): \texttt{yarn dev}
\end{enumerate}

The app then runs with a local database file in your application data folder.

\section{Main Navigation}
\label{sec:main-navigation}
The left sidebar has six main pages:
\begin{enumerate}
  \item \textbf{Dashboard}
  \item \textbf{Vendors}
  \item \textbf{Transactions}
  \item \textbf{Reports}
  \item \textbf{Goals}
  \item \textbf{Admin}
\end{enumerate}

\section{Dashboard}
The Dashboard gives a quick snapshot of the \textbf{most recent Saturday market day}.

\subsection{What You See}
\begin{itemize}
  \item Total reported sales
  \item Total reimbursement due
  \item Program token volume
  \item Sales by category chart
  \item Goal progress widget
  \item Vendor tracking table
\end{itemize}

\subsection{How to Use It}
\begin{itemize}
  \item Use it for a fast health check before or after market day.
  \item Search vendors in the table.
  \item Look for unusual alerts, like:
  \begin{itemize}
    \item an inactive vendor showing sales,
    \item or an active vendor with no sales entered.
  \end{itemize}
\end{itemize}

\section{Vendors}
Use this area to manage your master vendor list.

\subsection{What You Can Do}
\begin{itemize}
  \item View all vendors (active and optional inactive)
  \item Search by name, contact, email, location, or products
  \item Add a new vendor
  \item Edit an existing vendor
  \item Open vendor details
  \item Export vendor list to CSV
\end{itemize}

\subsection{Vendor Fields Explained}
\begin{longtable}{p{0.3\textwidth}p{0.65\textwidth}}
\textbf{Field} & \textbf{Meaning} \\
\hline
Vendor Name & Business name shown throughout the app. \\
Point Person & Main contact person. \\
Email & Contact email for vendor communication. \\
Location & Vendor location or stall reference. \\
Miles & Distance from market. \\
Products & What they sell (high-level and detailed notes). \\
Status & Active or Inactive. \\
Category Labels & Tags like Organic, Local, etc. \\
Defaults (percentages) & Product category split used in reports/estimates. \\
Average Sale Amount & Optional value used for transaction estimates. \\
\end{longtable}

\subsection{Add Vendor Workflow}
\begin{enumerate}
  \item Go to \textbf{Vendors} $\rightarrow$ \textbf{Add Vendor}.
  \item Fill in basic details.
  \item Add product category and product details.
  \item (Optional) Apply labels.
  \item (Optional) Set average sale amount.
  \item (Optional) Set category percentages.
  \item Review overview screen, then click \textbf{Create Vendor}.
\end{enumerate}

\textbf{Important rule:} if you enter category percentages, they must total exactly \textbf{100\%}.

\section{Transactions}
This is where you enter and maintain market-day transaction rows.

\subsection{Core Concept}
Each row is usually one vendor for one market date, including:
\begin{itemize}
  \item present/absent,
  \item token amounts (SNAP, DUFB, WDFM, voucher),
  \item reported sales,
  \item estimates and custom column values.
\end{itemize}

\subsection{Main Actions Menu}
On the Transactions page, use \textbf{Actions} for:
\begin{itemize}
  \item \textbf{Import Excel} --- upload worksheet data.
  \item \textbf{Download Template} --- get the official import format.
  \item \textbf{Add Vendor} --- add a vendor row manually.
  \item \textbf{Add Active Vendors} --- auto-add all currently active vendors.
\end{itemize}

\subsection{Typical Weekly Process}
\begin{enumerate}
  \item Select the correct market date.
  \item Add missing vendors to the sheet.
  \item Enter token and sales values.
  \item Fill required custom columns.
  \item Resolve any red/invalid rows.
  \item Save market-day statistics block (if used).
  \item Click \textbf{Save} to write all row changes.
\end{enumerate}

\subsection{Importing From Spreadsheet}
\begin{enumerate}
  \item Download template first.
  \item Fill it out.
  \item Click \textbf{Import Excel}.
\end{enumerate}

The app tries to match by transaction ID or vendor name and merges rows intelligently.

\subsection{Validation Checks Before Save}
The app blocks save when:
\begin{itemize}
  \item vendor names cannot be matched,
  \item required custom-column values are missing,
  \item percentage/default data is incomplete or invalid.
\end{itemize}

\section{Reports}
Reports are date-range based and can be exported as PDF.

\subsection{How Reports Work}
\begin{enumerate}
  \item Choose start and end date.
  \item Pick a report tab.
  \item Review charts/tables.
  \item Click \textbf{Export PDF}.
\end{enumerate}

\subsection{Report Tabs}
\begin{itemize}
  \item \textbf{Comprehensive}: overall totals, trend, and transaction table.
  \item \textbf{Category}: sales allocated by product category defaults.
  \item \textbf{Vendor Label}: sales split by vendor labels.
  \item \textbf{Leaderboard}: top vendors by total reported sales.
  \item \textbf{By Vendor}: deep dive for one vendor over time.
  \item \textbf{Token}: SNAP/DUFB/WDFM/voucher totals and related views.
  \item \textbf{Custom Columns}: rollups for admin-defined transaction fields.
\end{itemize}

\subsection{Reading Report Numbers}
\begin{itemize}
  \item \textbf{Reported Sales}: what vendors reported for the day/range.
  \item \textbf{Reimbursement Due}: amount owed back based on token redemptions.
  \item \textbf{Token Volume}: combined token-related amounts.
\end{itemize}

\section{Goals}
Goals let you track targets over a date range.

\subsection{What You Can Do}
\begin{itemize}
  \item Create goal
  \item Edit goal
  \item Delete goal
\end{itemize}

\subsection{Goal Setup}
Each goal includes:
\begin{itemize}
  \item Name
  \item Start date and end date
  \item Metric (for example, reported sales or attendance percentage)
  \item Target value
\end{itemize}

The app automatically calculates:
\begin{itemize}
  \item current value from transaction data,
  \item percent toward goal.
\end{itemize}

\textbf{Note:} Attendance-type goals are limited to \textbf{0--100\%}.

\section{Admin}
Admin is for configuration and system-level maintenance.

\subsection{Category Labels}
\begin{itemize}
  \item Create, edit, delete labels.
  \item Labels can be color-coded.
  \item Labels are reusable across vendors.
\end{itemize}

\subsection{Custom Columns}
Add fields that appear on transaction rows for your whole market.

\textbf{Column types:}
\begin{itemize}
  \item text
  \item number
  \item boolean (yes/no)
  \item usd (currency)
\end{itemize}

You can also:
\begin{itemize}
  \item mark fields required,
  \item deactivate/reactivate columns,
  \item permanently delete inactive columns.
\end{itemize}

\subsection{Database Management (Desktop Mode)}
\begin{itemize}
  \item \textbf{Export Database}: create a backup \texttt{.mv.db} file.
  \item \textbf{Import Database}: restore from a backup file.
\end{itemize}

\textbf{Warning:} database import overwrites current data.

\section{Data and Safety Tips}
\subsection{Before Market Day}
\begin{itemize}
  \item Verify active vendor list is current.
  \item Verify required custom columns are set correctly.
  \item Download fresh transaction template if importing.
\end{itemize}

\subsection{After Market Day}
\begin{itemize}
  \item Save transactions.
  \item Review dashboard totals.
  \item Run reports for archive and stakeholders.
  \item Export database backup.
\end{itemize}

\section{Common Problems and Fixes}
\subsection{``No data'' in reports}
\begin{itemize}
  \item Check date range.
  \item Confirm transactions were saved for those dates.
\end{itemize}

\subsection{Cannot save transactions}
\begin{itemize}
  \item Fix unmatched vendor rows.
  \item Fill required custom fields.
  \item Ensure percentage/default requirements are valid.
\end{itemize}

\subsection{Import did not map correctly}
\begin{itemize}
  \item Use official template.
  \item Keep vendor names consistent with vendor master list.
\end{itemize}

\subsection{Database backup/import buttons do not work}
Those actions are available only in the desktop app runtime (Tauri), not standard browser mode.

\Needspace{12\baselineskip}
\section{Glossary}
\begin{tabularx}{\textwidth}{p{0.3\textwidth}X}
\textbf{Term} & \textbf{Meaning} \\
\hline
Vendor Defaults & Saved category percentages and average sale estimate for a vendor. \\
Market Day Data & Market-level token stats for a specific date. \\
Reimbursement Due & Amount to be paid back to vendor based on eligible tokens. \\
Custom Column & Admin-defined transaction field used by all vendors. \\
Label & Reusable tag assigned to vendors for grouping/reporting. \\
\end{tabularx}

\section{Where This Manual Lives}
This manual file is stored in the repository at:
\begin{itemize}
  \item \texttt{docs/user-manual.tex}
\end{itemize}

To produce a PDF from this file, run a standard LaTeX build command (for example, \texttt{pdflatex user-manual.tex}) in the \texttt{docs} folder.

\end{document}
